\chapter{Un socle pour la suite}
\label{ch:socle}

\questionchapitre{Ces développements survivront-ils au projet qui les a financés, et
remplissent-ils la mission qui leur était assignée, à savoir outiller les travaux
d'explicabilité à venir ?}

\section{La question posée franchement}

Un logiciel de recherche est ordinairement écrit pour produire un résultat, puis abandonné. Il
suppose un environnement particulier, mêle le traitement et son affichage, dépend de fichiers
qui ne sont plus là, et sa reprise coûte plus cher que sa réécriture.

Affirmer qu'un développement constitue un socle n'a donc de valeur que si l'on peut désigner les
propriétés qui le rendent réutilisable et montrer qu'elles sont acquises. C'est l'objet de ce
chapitre. La question est d'autant plus directe que la mission était explicitement de
\emph{préparer} des travaux ultérieurs : un socle qui ne serait pas repris constituerait un
échec de la mission, indépendamment de sa qualité technique.

\section{Les propriétés acquises}

\subsection{Un contrat d'interface qui découple l'amont de l'aval}

La couche analytique de l'entrepôt se consomme en SQL. Cette interface est stable, universelle
et documentée : un utilisateur n'a besoin de connaître ni les outils d'ingestion, ni les
bibliothèques de transformation, ni l'orchestrateur. Un carnet de calcul, un tableur connecté,
un outil de visualisation ou une application tierce s'y branchent également.

La conséquence est structurante : l'amont peut être remplacé sans que l'aval s'en aperçoive. Si
un outil d'ingestion devient obsolète, il est substituable sans affecter les consommateurs. Peu
de logiciels de recherche offrent cette propriété, et c'est elle qui distingue une
infrastructure d'un script.

\subsection{Un c\oe ur de calcul indépendant de toute interface}

Le c\oe ur métier de la plateforme ne comporte aucune dépendance à l'interface graphique
(chapitre~\ref{ch:plateforme}). Le même code est appelé depuis le service web, un script ou un
carnet de calcul. Même si l'application était abandonnée, les méthodes resteraient utilisables.
La contribution scientifique n'est pas prisonnière de son emballage.

\subsection{Des indicateurs validés une fois pour tous les usages}

Les indicateurs du chapitre~\ref{ch:indicateurs} sont calculés dans l'entrepôt, validés selon un
protocole documenté, et lus tels quels par tous les consommateurs. Il n'existe pas de second
calcul susceptible de diverger. Cette centralisation garantit que deux études conduites à des
dates différentes, ou par deux équipes, parlent de la même chose, condition de comparabilité que
l'on ne peut pas rattraper après coup.

\subsection{Des contrats exécutables entre composants}

Là où une duplication était nécessaire, la classification en classes de sévérité présente dans
les deux dépôts, elle est gardée par des tables de référence identiques et des tests qui
échouent en cas de divergence. La contrainte est exécutable plutôt que documentaire : elle
résiste à l'oubli, ce qu'aucune note ne fait.

\subsection{Une infrastructure reproductible et éprouvée}

L'ensemble est conteneurisé, les volumes de données sont déclarés hors du cycle de vie des
conteneurs, et le chargement initial est interruptible et redémarrable. Les procédures de
sauvegarde et de restauration ont été vérifiées à l'échelle réelle. Une réinstallation complète
est une procédure documentée et testée, non une reconstitution.

\subsection{Une documentation écrite pour la reprise}

La documentation des deux dépôts a été refondue en fin de période dans une perspective explicite
de passation : cartographie des documents, procédures vérifiées par exécution réelle, et
consignation des pièges rencontrés, y compris ceux qui produisent des résultats faux sans
émettre d'erreur. Ce dernier point est le plus utile à un successeur et le plus rarement écrit.

\section{Le patron généralisable}

Au-delà des deux dépôts, ces travaux instancient un enchaînement (figure~\ref{fig:patron}) qui
n'a rien de spécifique aux eaux souterraines.

\begin{figure}[H]
\centering
\begin{tikzpicture}[
  font=\sffamily\small,
  et/.style={draw=black!70, fill=black!3, minimum width=2.6cm, minimum height=1.3cm,
             align=center, font=\sffamily\scriptsize},
  fl/.style={-{Stealth[length=2.4mm]}, draw=black!70, thick},
  note/.style={font=\sffamily\scriptsize\itshape, color=black!55, align=center,
               text width=2.55cm},
]
\node[et] (e1) {\textbf{Sources}\\ouvertes,\\hétérogènes};
\node[et, right=0.7cm of e1] (e2) {\textbf{Entrepôt}\\consolidé,\\interrogeable};
\node[et, right=0.7cm of e2] (e3) {\textbf{Indicateurs}\\normalisés\\et validés};
\node[et, right=0.7cm of e3] (e4) {\textbf{Modèles}\\métier et\\appris};
\node[et, right=0.7cm of e4, very thick, draw=black] (e5) {\textbf{Explication}\\et\\scénarisation};

\foreach \a/\b in {e1/e2, e2/e3, e3/e4, e4/e5} \draw[fl] (\a) -- (\b);

\node[note, below=0.32cm of e2] {coût payé une seule fois};
\node[note, below=0.32cm of e3] {comparabilité garantie ici};
\node[note, below=0.32cm of e5] {condition de l'usage en décision};
\end{tikzpicture}
\caption{L'enchaînement mis en \oe uvre. Chaque étape suppose la précédente ; escamoter l'une
d'elles produit des résultats incomparables ou inutilisables.}
\label{fig:patron}
\end{figure}

L'enseignement de la période tient dans la lecture de gauche à droite de cette figure : on ne
peut pas commencer par la droite. C'est pourtant ce que font de nombreux projets d'intelligence
artificielle appliquée à l'environnement, qui investissent dans le modèle et découvrent ensuite
que la donnée ne permet pas de conclure.

\section{Ce qui est transmis aux travaux d'explicabilité}

La mission consistait à préparer le terrain d'\textsc{Aida}. Le tableau suivant confronte les
besoins de ce projet à ce qui lui est effectivement transmis.

\begin{center}
\begin{tabularx}{\textwidth}{@{}p{4.6cm}X@{}}
\toprule
\textbf{Besoin} & \textbf{Disponible} \\
\midrule
Un jeu de données consolidé, joint au forçage climatique, sur lequel entraîner des modèles &
Entrepôt national, historique depuis 1967 pour les stations et 1950 pour le climat,
rattachement spatial effectué et matérialisé \\
Des modèles à expliquer & Douze architectures de prévision profonde sous interface unique, avec
traçabilité des expériences \\
Une référence interprétable à laquelle confronter les explications & Modèles à fonction de
transfert calibrés, dont les paramètres ont un sens hydrogéologique \\
Des méthodes d'attribution & Cinq familles complémentaires, confrontables entre elles \\
Des explications actionnables & Contrefactuels sous contrainte physique, mettant en \oe uvre les
critères formulés par les hydrogéologues \\
Un moyen de valider une explication & Rejeu du scénario dans un modèle métier indépendant \\
Des cas d'étude qualifiés & Détection non supervisée d'influence anthropique, permettant
d'écarter les chroniques impropres \\
\bottomrule
\end{tabularx}
\end{center}

Ce qui n'est pas transmis doit être dit avec la même netteté : aucun résultat de performance
comparée, aucune évaluation quantitative des méthodes contrefactuelles, et aucune campagne
d'expériences consolidée. Le socle rend ces travaux possibles ; il ne les a pas menés.

\section{Extensions, par coût croissant}

\begin{center}
\begin{tabularx}{\textwidth}{@{}lXl@{}}
\toprule
\textbf{Extension} & \textbf{Ce qu'elle suppose} & \textbf{Coût} \\
\midrule
Nouvelles stations, autre territoire & Rien : la couverture est nationale et la chaîne est
paramétrée par configuration & nul \\
Nouveaux indicateurs standardisés & Un ajustement statistique et sa validation, selon le
protocole existant & faible \\
Nouvelles architectures de prévision & Une entrée au catalogue ; la fabrique de modèles et la
traçabilité sont génériques & faible \\
Qualité des eaux souterraines & De nouvelles sources et de nouveaux modèles de transformation ;
l'architecture est inchangée & moyen \\
Eaux de surface & La chaîne hydrométrique existe déjà ; l'extension porte sur la modélisation &
moyen \\
Autre compartiment environnemental & Les sources et les indicateurs diffèrent, le patron tient &
moyen \\
Couplage à un modèle physique distribué & Une interface avec un simulateur externe et
l'assimilation de ses sorties & élevé \\
\bottomrule
\end{tabularx}
\end{center}

Le tableau appelle une lecture précise : ce qui est peu coûteux l'est \emph{parce que}
l'infrastructure existe. Le coût élevé de l'année écoulée a été payé une fois et ne se
représente pas à chaque extension. C'est la définition d'un socle.

\section{Ce qui manque pour un environnement d'outils partagé}

Cinq chantiers séparent l'état actuel d'un environnement partagé. Aucun n'est un verrou de
recherche ; tous sont des décisions d'investissement.

\begin{enumerate}
  \item \textbf{L'automatisation des vérifications.} La suite de tests de la plateforme n'est
  pas exécutée automatiquement (chapitre~\ref{ch:apprentissage}). Tant que les vérifications
  dépendent d'une initiative humaine, la confiance dans les évolutions décroît avec le temps.
  C'est le chantier le moins coûteux et le plus rentable.

  \item \textbf{Un accès élargi et gouverné.} L'accès nominatif convient à un outil de projet,
  non à un environnement partagé. Une politique d'accès conditionne tout élargissement.

  \item \textbf{Un catalogue de données.} L'entrepôt est documenté pour qui sait où chercher. Un
  environnement d'outils suppose un catalogue décrivant les jeux disponibles, leur origine, leur
  fraîcheur et leurs conditions d'usage.

  \item \textbf{Le calcul mutualisé.} L'entraînement s'appuie sur un serveur unique. Un usage
  partagé suppose une file d'attente et un partage de ressources.

  \item \textbf{La campagne d'évaluation comparative.} Le socle sera d'autant plus crédible
  qu'il aura servi à établir un résultat. L'outillage est prêt de part et d'autre.
\end{enumerate}

\begin{acquisouvert}
\textbf{Acquis.} Les propriétés qui distinguent une infrastructure d'un prototype sont
identifiées et tenues : interface stable, c\oe ur indépendant, indicateurs centralisés, contrats
exécutables, reproductibilité, documentation de reprise. Le patron mis en \oe uvre est
transposable à d'autres compartiments environnementaux, et les besoins des travaux
d'explicabilité à venir sont couverts point par point.

\textbf{Ouvert.} Cinq chantiers séparent ce socle d'un environnement partagé, et la
démonstration par un résultat scientifique reste à produire.
\end{acquisouvert}
